% !TeX root = ../main.tex


\section{Tecnologías usadas}
\label{sec:tecnologias_usadas}
Debido a la naturaleza computacional de este proyecto, el lenguaje principal de programación empleado ha sido \texttt{Python}, destacando por su legibilidad y amplia adopción en la comunidad astrofísica. Como base para la manipulación y vectorización de datos se utilizó \texttt{Numpy}, mientras que la librería \texttt{Astropy} facilitó el manejo de datos observacionales reales en formato \texttt{FITS}. 

Para el procesamiento de las imágenes, se empleó \texttt{Scipy}, cuyos módulos proporcionan herramientas esenciales como la convolución, transformadas de Fourier y algoritmos de optimización (aplicados en las secciones \ref{sec:metodosDeconvolucion} y \ref{sec:metodoDes}). Además, esta librería permitió paralelizar fácilmente el código.

Por último, se ha utilizado la inteligencia artificial como una potente herramienta de apoyo metodológico durante todo el proyecto, destacando en tres áreas principales:

\begin{itemize}
    \item Para facilitar y acelerar la comprensión de conceptos teóricos complejos abordados en la literatura científica.
    \item Para agilizar la localización de errores en el código y asistir en la generación de scripts para la visualización de datos y gráficas.
    \item Como soporte ortográfico y de corrección de estilo para mejorar la calidad en la redacción de este documento.
\end{itemize}